\section{The game measures the embezzlement error}\label{sec:equiv}

We first collect the properties of $\cE$ that are used below. The
bounds~\eqref{eq:window}--\eqref{eq:variance} are contained in Ref.~\cite[Section~6.4]{TG}, and
we include short proofs to keep the paper self-contained. Throughout, $T$ denotes a random
variable on $(0,\infty)$ with density $p(t)=p_{\lceil t\rceil}$, which is a continuous version
of the rank of a Schmidt probability, and $X=\log_2T$.

\begin{proposition}[Properties of $\cE$]\label{prop:E}
Let $\psi$ be a pure bipartite state with Schmidt probabilities $p_1\ge p_2\ge\cdots$, Schmidt
coefficients $\lambda$ and entanglement entropy $S=-\sum_jp_j\ln p_j\le\infty$. The embezzlement
error is positive and never exceeds its value on a product state,
\begin{equation}\label{eq:range}
0<\cE(\psi)\le2-\sqrt2,
\end{equation}
with equality on the right exactly when all nonzero $p_j$ are equal, that is, for product
states and maximally entangled states. It does not change when a maximally entangled state
$\Phi_k$ of any local dimension $k$ is added, and its square root is Lipschitz continuous in the
Schmidt coefficients: if $\psi'$ has Schmidt coefficients $\lambda'$, then
\begin{equation}\label{eq:invariance}
\cE(\psi\otimes\Phi_k)=\cE(\psi),\qquad
\bigl\lvert\sqrt{\cE(\psi)}-\sqrt{\cE(\psi')}\bigr\rvert\le2\norm{\lambda-\lambda'}.
\end{equation}
A small error forces the Schmidt probabilities to spread over many ranks. The largest of them,
and the total weight of the ranks between $r+1$ and $2r$ for every integer $r\ge1$, satisfy
\begin{equation}\label{eq:window}
p_1\le(6+4\sqrt2)\,\cE(\psi),\qquad\sum_{j=r+1}^{2r}p_j\le\sqrt{\cE(\psi)} ,
\end{equation}
while the entanglement entropy and, when finite, the variance of $X$ satisfy
\begin{equation}\label{eq:entropybound}
\sqrt{\cE(\psi)}\;\ge\;\frac{\ln2}{2\,(S+1+3\ln2)}
\end{equation}
and
\begin{equation}\label{eq:variance}
\cE(\psi)\,\bigl(4\Var(X)+1\bigr)\ge1 .
\end{equation}
\end{proposition}

\begin{proof}
We use the formula $\cE=\norm{\lambda-D\lambda}^2=2-2\ip{\lambda}{D\lambda}$ of
Lemma~\ref{lem:formula} throughout.

\emph{Range.} Since $\lambda$ is decreasing, $\lambda_{\lceil m/2\rceil}\ge\lambda_m$. Hence
$\ip{\lambda}{D\lambda}=\sum_m\lambda_m\lambda_{\lceil m/2\rceil}/\sqrt2\ge1/\sqrt2$, which is
the upper bound in~\eqref{eq:range}. Equality requires
$\lambda_m\in\{0,\lambda_{\lceil m/2\rceil}\}$ for all $m$, that is, a flat spectrum. Positivity
follows from the first bound in~\eqref{eq:window}.

\emph{Invariance and continuity.} The Schmidt coefficients of $\psi\otimes\Phi_k$ are
$(D_k\lambda)_m=\lambda_{\lceil m/k\rceil}/\sqrt k$. The map $D_k$ is an isometry, and it
commutes with $D=D_2$ because $\lceil\lceil m/k\rceil/2\rceil=\lceil m/(2k)\rceil$. Therefore
$\norm{D_k\lambda-DD_k\lambda}=\norm{D_k(\lambda-D\lambda)}=\norm{\lambda-D\lambda}$. The
Lipschitz bound holds because $\sqrt{\cE}=\norm{(1-D)\lambda}$ and $\norm{1-D}\le2$.

\emph{Largest probability and windows.} The first component of $\lambda-D\lambda$ is
$\lambda_1(1-1/\sqrt2)$, so $\cE\ge p_1(1-1/\sqrt2)^2=p_1/(6+4\sqrt2)$. For the second bound
in~\eqref{eq:window}, let $f(s)=(\ln2)\,2^sp(2^s)$ be the density of $X$ and $g=\sqrt f$. The
substitution $t=2^{s+1}$ in
$\norm{\lambda-D\lambda}^2=\int_0^\infty\bigl(\sqrt{p(t)}-\sqrt{p(t/2)/2}\bigr)^2dt$ gives
\begin{equation}\label{eq:Eshift}
\cE(\psi)=\int_{\R}\bigl(g(s+1)-g(s)\bigr)^2\,ds .
\end{equation}
In other words, the embezzlement error measures how much the distribution of $X$ changes when
it is shifted by one. Put $\mu(u)=\int_u^{u+1}f$. By the Cauchy--Schwarz inequality,
$\sum_{k\in\Z}\lvert\mu(u+k)-\mu(u+k-1)\rvert\le\int_{\R}\lvert g(s)-g(s-1)\rvert\,(g(s)+g(s-1))\,ds\le2\sqrt{\cE}$.
The left-hand side is the total variation of a nonnegative summable sequence and is therefore
at least twice its maximum. So $\mu(u)\le\sqrt{\cE}$ for all $u$, and the choice $u=\log_2r$
gives the claim.

\emph{Entropy.} Because $p$ is decreasing, $t\,p(t)\le\int_0^tp\le1$, hence
$-\ln p(T)\ge\ln T$ and $S\ge(\ln2)\,\mathbb{E}[X]$. The negative part of $X$ contributes
$\int_0^1p_1\log_2(1/t)\,dt=p_1/\ln2\le1/\ln 2$. Put $h=\sqrt{\cE}$. For $x\ge0$ the interval
$(1,2^x]$ is covered by at most $x+1$ windows of the form $(r,2r]$, and $p_1=2\mu(-1)\le2h$, so
$\Pr[X\le x]\le(x+3)h$ and
$\mathbb{E}[X_+]\ge\int_0^\infty(1-(x+3)h)_+\,dx=(1-3h)_+^2/(2h)$. If $h<1/3$ this gives
$S\ge(\ln2)\bigl(\tfrac1{2h}-3\bigr)-1$, which is~\eqref{eq:entropybound}. If $h\ge1/3$
then~\eqref{eq:entropybound} holds because its right-hand side is below $0.12$.

\emph{Variance.} Assume $X$ has finite second moment; otherwise the variance bound is interpreted as vacuous. The laws of $X$ and $X-1$ have means that differ by one. With
$\mu_X=\mathbb{E}[X]$ and $c=\mu_X-\tfrac12$ we obtain
$1=\int(s-c)\bigl(f(s)-f(s+1)\bigr)ds\le\sqrt{\cE}\,\bigl(\norm{(s-c)g(s)}_2+\norm{(s-c)g(s+1)}_2\bigr)
=2\sqrt{\cE}\,\sqrt{\Var(X)+\tfrac14}$, which is~\eqref{eq:variance}.
\end{proof}

\begin{remark}\label{rem:sharp}
The power of $S$ in~\eqref{eq:entropybound} cannot be improved for general states, since the sine profiles in Section~\ref{sec:operational-lower} have
$\cE=O(K^{-2})$ and $S=\Theta(K)$. Their shell distribution keeps a fixed
fraction of its mass at ranks $2^{\Theta(K)}$, giving the entropy estimate. For exact conformal tower spectra $\cE$
is of order $1/S$ (Section~\ref{sec:law}), and~\eqref{eq:variance} is then the relevant bound.
The second bound in~\eqref{eq:window} controls the weight of every factor-of-two window of
ranks. In particular the quantity $\eta_{\mathrm{emb}}$ of Ref.~\cite{Sierant26}, which decides
universal LOCC embezzlement~\cite{ZTG24}, satisfies $\eta_{\mathrm{emb}}\le2\sqrt{\cE}$. Explicit
Gaussian embezzlement protocols for free-fermion ground states are given in
Ref.~\cite{KvLSW25}, which also notes that their convergence with the system size is very slow.
\end{remark}


\paragraph{The numerical lower comparison.}
The proof in Appendix~\ref{sec:explicit-robust} extracts the magic-square
sector with a $210\sqrt N$ error, where $N$ is its total forbidden mass.
A polar decomposition then aligns the nine label sectors, at a cost bounded
by twice the square root of their total mismatch. These estimates give
\[
 \mathcal E(\psi)\le C\,\varepsilon,\qquad
 C=\frac{243(210^2+448)}8=\frac{2{,}706{,}291}{2},
\]
for every strategy of loss $\varepsilon$. Local projective dilations preserve
the nonzero Schmidt spectrum, so the bound holds for arbitrary POVMs and
arbitrary Schmidt rank. Taking the infimum of the loss proves the lower
bound in Theorem~\ref{thm:equiv}.
